\section{Introduction}
\label{sec:intro}

\noindent
Large language models (LLMs) have become an important cloud service
thanks to their emerging capabilities in various tasks.
Both customer-facing services (to-C) like conversation-based chatbots~\cite{chatgpt}
and business services (to-B) through LLM API calls have been widely deployed on the cloud~\cite{openaiapi,genimiapi,claudeapi}.
As online services, serving LLMs with low latency and high throughput is critical,
but achieving this in a cost-efficient way is challenging
due to the huge computational requirements to serve each LLM request.

Caching intermediate results of served requests ({\kvcache}) has become a common methodology
to improve serving performance~\cite{openai-cache,claude-cache}:
when two requests share the same input prefix,
their {\kvcache} are identical;
so if we cache the {\kvcache} of the first ({\kvcache} cache),
we don't need to compute the {\kvcache} for the next.
Hence, the serving latency is reduced with the throughput improved
thanks to the reduced computation under hits.

Like traditional data caching systems~\cite{yang2020twemcache,atikoglu2012workload,180324},
the effectiveness of {\kvcache} cache system components 
like cache eviction policy should be tailored to the {\kvcache} reuse features of workloads.
However,
LLM serving workloads are complex due to their diverse request types
as well as the unique features of each request type.
For example, a production LLM service handles both to-C workloads (like chatting)
and to-B workloads (like task automation with API calling).
These requests can further be processed either as single-turn or multi-turn interactions (\textsection{\ref{sec:bg}}).
Since the {\kvcache} can be reused across different single-turn requests and multi-turn requests,
both the single-turn request types and the multi-turn patterns
can possibly affect the design choices of {\kvcache} cache.

There remains a significant gap in the understanding of common LLM serving
workloads for {\kvcache} cache due to the lack of in-depth analysis of real-world traces.
First, it is unclear whether production workloads have many reuses and which request
type contributes most to the reuses.
Second, the distributions of the reuse time or reuse probability of {\kvcache} remain unknown,
which is critical to cache policy designs.
Finally, there is a lack of characterization of the lifespan of {\kvcache},
which is essential in helping us estimate the required cache capacity for serving LLMs in the wild.
\iffalse
are critical to the effectiveness of cache policies.
\emph{Third}, how long does the {\kvcache} cache last?
The cached {\kvcache} can become obsolete quickly rather than being reusable long-term, which differs from traditional caching.
The lifespan of the {\kvcache} cache helps estimate the {\kvcache} capacity for the serving workloads and guides the design of a {\kvcache} system such as which storage media to use.
\fi

\iffalse
First, there is a lack of comprehensive studies examining whether {\kvcache} caching
is effective in production LLM serving workloads and which request types contribute most to cache hits.
Second, whether a {\kvcache} block will be hit in the future depends on the workload patterns,
such as the probability of the reuse and the interval between the next reuse.
Such patterns are critical to the effectiveness of cache policies,
but to the best of our knowledge---are not revealed on large-scale
production traces.
Finally, {\kvcache} differs from traditional caching in its lifespan---cached data
can become obsolete quickly rather than being reusable long-term.
Understanding this temporal property helps estimate the {\kvcache} capacity for the serving workloads,
which further guides the design of a {\kvcache} system like which storage media to use.
Such information is still missing in existing studies.
\fi


\begin{table}[!t]
    \centering
%    \vspace{2.2mm}
    \small{
        \resizebox{1\linewidth}{!}{
            \ra{1.2}
            \begin{tabular}{lrrrrr}
                \toprule
                & \multicolumn{1}{c}{\textbf{Time}} & \multicolumn{1}{c}{\textbf{Type}} & \multicolumn{1}{c}{\textbf{UID}} & \multicolumn{1}{c}{\textbf{Turn}}     & \multicolumn{1}{c}{\textbf{Content}} \\ \hline
\textbf{ShareGPT~\cite{sharegpt}}  &      \xmark                    &   \xmark                    &        \xmark           & \cmark                                  &       \cmark                               \\
\textbf{Mooncake~\cite{mooncake}}  &      \cmark                    &   \xmark                    &        \xmark           & \xmark                                  &       \cmark                               \\
\textbf{Traces in this work}       &      \cmark                    &   \cmark                    &        \cmark           & \cmark                                  &       \cmark                               \\ 
                \bottomrule
\end{tabular}
            
        }
    } \\[8pt]
    \begin{minipage}{1\linewidth}
        \caption{\small{\emph{
        A comparison of currently available LLM serving traces. \textbf{Time} is the invocation time of each request. 
        \textbf{Type} is the request type, e.g., whether it is an API call. 
        \textbf{UID} is the user ID that invokes the request. 
        \textbf{Turn} is the multi-turn information of the request. 
        \textbf{Content} is the (anonymous) input and output tokens of each request.
        }}}
    \label{tab:compare-traces}
    \end{minipage} \\[-15pt]
\end{table}  


% \begin{table}[!t]
%     \centering
%     \vspace{2.2mm}
%     \small{
%         \resizebox{.99\linewidth}{!}{
%             \ra{1.2}

%             \begin{tabular}{lrrrrr}
%                 \toprule
%                 & \multicolumn{1}{c}{\textbf{Time}} & \multicolumn{1}{c}{\textbf{Type}} & \multicolumn{1}{c}{\textbf{UID}} & \multicolumn{1}{c}{\textbf{Turn}}     & \multicolumn{1}{c}{\textbf{Content}} \\ \hline
% \textbf{ShareGPT~\cite{sharegpt}}  &      \xmark                    &   \xmark                    &        \xmark           & \cmark                                  &       \cmark                               \\
% \textbf{Mooncake~\cite{mooncake}}  &      \cmark                    &   \xmark                    &        \xmark           & \xmark                                  &       \cmark                               \\
% \textbf{Our trace}                 &      \cmark                    &   \cmark                    &        \cmark           & \cmark                                  &       \cmark                               \\ 
%                 \bottomrule
% \end{tabular}
            
%         }
%     } \\[8pt]
%     \begin{minipage}{1\linewidth}
%         \caption{\small{\emph{
%         A comparison of prefix cache works.
%         }}}
%     \label{tab:compare-works}
%     \end{minipage} \\[-10pt]
% \end{table}  


\stitle{Characterizing production traces. \,}
To bridge this gap,
we collected and characterized two representative production serving workloads
from one of the world's largest cloud providers (\company):
a to-C workload where users submit requests to the cloud-hosted LLM service via browser-based
chatbots or mobile applications, and a to-B workload where developers
send OpenAI-compatible requests to the LLM service through API calls in their programs.
We collected as much detailed information as possible within the company's strict privacy policy,
including request submission times, request types (single-turn, multi-turn, API requests, or others),
and anonymized content to analyze caching effectiveness.
To our knowledge, no other available LLM serving traces
match the information depth and scale of our traces (summarized in Table~\ref{tab:compare-traces}).
For example, ShareGPT~\cite{sharegpt} only provides the input and output of each request,
not the submission time of each request,
which is critical for analyzing the real {\kvcache} hit ratio and usage over time.
We believe our collected workloads represent typical LLM serving workloads
and can provide a foundation for future {\kvcache} cache system design.

From the perspective of designing an efficient {\kvcache} cache system---both in performance
and resource usage,
our \emph{key takeaways} are as follows: \\[-15pt]
\begin{enumerate}[leftmargin=*,leftmargin=10pt,itemindent=0pt]
    \item {\kvcache} reuses are common,
    but the reuse ratio is smaller than previously reported numbers
    on synthetic datasets~\cite{cachedattention,DBLP:journals/corr/abs-2403-14401}.
    The reuses follow a skewed distribution:
    10\,\% of {\kvcache} blocks contribute to 77\,\% of the reuses.
    The contribution of {\kvcache} reuses also varies across traces:
    contrary to a straightforward observation that multi-turn requests may dominate most reuses,
    single-turn requests dominate 97\,\%
    {\kvcache} reuses in to-B workloads (\textsection{\ref{sec:workload-patterns}}).  \\[-15pt]

    \item Key workload features related to {\kvcache} designs,
    like reuse time and reuse probability distributions---vary across different LLM request types (API vs. chat)
    and the number of turns in multi-turn requests.
    Despite this diversity, for each specific request category (type plus turn number),
    the reuse time is predictable based on the historical information (\textsection{\ref{sec:motiv-temporal-spatial}}).
    Since the request categories are known by the serving system,
    we can systematically use the estimated probability for improving cache policies. \\[-15pt]

    \item The lifespan of {\kvcache} is ephemeral: the P99 lifespan
    of {\kvcache} in to-B workloads is 97 seconds,
    so a small cache is typically sufficient for the {\kvcache} cache.
    For example, in the to-B workload,
    a {\kvcache} with capacity 2\,$\times$ of the GPU HBM per-GPU
    is sufficient to approach an ideal hit rate under infinite capacity on
    common GQA models~\cite{DBLP:conf/emnlp/AinslieLJZLS23}
    using standard cache eviction policies like LRU (\textsection{\ref{sec:analyze-capacity}}).
\end{enumerate}

Our characterizations above reveal several important design decisions
that current {\kvcache} cache systems may have overlooked.
First, optimizing caching for single-turn requests is equally important as optimizing
for multi-turn requests.
Second, we could improve cache hits
by adopting a workload-aware caching policy
instead of only adopting a workload-agnostic policy like LRU.
Finally, the ephemeral lifespan of {\kvcache} suggests that
a least-frequently-used (LFU) policy is not suitable because a short-lived {\kvcache} 
may have high frequency in the past. 
Meanwhile, 
a small (on-GPU) cache is sufficient for API-dominated to-B workloads,
eliminating the cost and complexity of
deploying and managing a CPU-RDMA-SSD storage hierarchy for {\kvcache} cache.
Besides the main takeaways described above, 
we also discovered several interesting findings, 
such as the {\kvcache} hits across users being extremely low (\textsection{\ref{sec:workload-patterns}}),
despite the possibility of shared prompts via prompt libraries~\cite{prompt-library}. 
%and unlike reuse time, the distribution of lifespan for each {\kvcache} being hard to predict (\textsection{\ref{sec:analyze-capacity}}). 

\stitle{Improved {\kvcache} cache eviction policy. \,}
Based on our findings,
we propose a workload-aware {\kvcache} eviction policy (\textsection{\ref{sec:design-policy}}) that leverages
the profiled reuse probability distributions of each workload to enhance cache hit ratios
on real-world workloads.
%Specifically, we built upon the Greedy-Dual caching framework~\cite{cherkasova1998improving}---which
%already considers various workload features---by further extending
Specifically, we leverage the empirically sampled probability distributions of the reuse probability 
of each workload to decide the priority of each {\kvcache} block, 
with considerations of both the characterized spatial locality and the lifespan of {\kvcache}.
To demonstrate the effectiveness of our improved design,
we have integrated our policy into vLLM~\cite{vllm}---a popular LLM serving system on the cloud,
and our extensive evaluations on production traces show that
the improved design can improve 3.9\,\% cache hits
as well as up to 41.4\,\% mean response time improvements
compared to workload-agnostic policies like LRU and LFU.
%It also has 3.8\,\% cache hit rate improvement and 34.1\,\% mean response time reduction
%over the original Greedy-Dual policy.

\stitle{Discussion. \,}
Before moving on, we note two limitations of our work.
First, our results are based on one week production traces of large-scale LLM serving workloads.
While we believe such workloads (chatbot and API) are representative nowadays,
LLM serving workloads are evolving, with new workloads like reasoning~\cite{openai-o1} emerging.
We leave the characterization of these newer workloads to future work.
Second, our primary focus is on the {\kvcache} cache policy design,
yet we believe the design of other system components for LLM serving 
like the policy of global scheduling 
could also benefit from our findings,
and we leave other explorations as future work.
Nevertheless, we believe our methodology---first characterizing real-world serving features
and then optimizing a specific serving component---generalizes to other LLM serving workloads and systems.

%We may open-source our traces if approved by the company.
A sample of our analyzed trace is available at \burl{https://github.com/alibaba-edu/qwen-bailian-usagetraces-anon},
and the reference implementation of the eviction policy is available at
\burl{https://github.com/vllm-project/vllm/pull/22236}. 